\section{Conclusion}\label{sec:conclusion}

Architecture choice has a modest effect on peptide-MHC binding prediction: a CNN ranked above the Deep FFN in all five seeds; the top four architectures fell within $\sim$2 percentage points, and a Random Forest was clearly worse. Random-label controls confirmed chance-level performance, ruling out noise overfitting. Canonical binding is largely captured by the anchor motif, but a zero-parameter anchor rule (sensitivity 0.79) underperforms the trained model (0.98), which recovers non-canonical binders the rule misses --- while non-canonical binders reach $\sim$0.65 sensitivity across four sequence-based predictors at unmatched operating points, marking the unresolved part of the problem. Beyond peptide-MHC, the three artifacts we corrected --- data-size imbalance, a degenerate null, and stale checkpoints --- illustrate a general benchmark-composition failure mode that reproduces false conclusions across machine-learning benchmarks; the checklist in Box~1 is offered as a reusable safeguard. Experimental validation of the non-canonical tail is the decisive next step.
